\documentclass[12pt,a4paper]{article}
\input{preamble}

\renewcommand{\ReportKind}{Assignment 1}
\renewcommand{\ReportSubtitle}{Nền tảng đặt cà phê trực tuyến\\ thanh toán không dùng tiền mặt}
\renewcommand{\RunningHead}{BrewLite -- Assignment 1: Use case, User story \& Kiến trúc}

% --- TikZ ERD styles -------------------------------------------------
\tikzset{
  ent/.style={draw=coffee,fill=cream,rounded corners=2pt,inner sep=4.5pt,
              font=\scriptsize,align=left,text width=2.95cm},
  lnk/.style={draw=coffee!85,thick},
  card/.style={font=\scriptsize\bfseries,text=coffee,fill=white,inner sep=1pt},
  note/.style={font=\scriptsize\itshape,text=coffee!85,align=center},
}

\begin{document}
\makecover

\section{Giới thiệu}

\subsection{Bối cảnh và vấn đề}
BrewLite là kênh đặt cà phê trực tuyến cho một quán quy mô nhỏ. Trước khi có hệ thống, quán nhận
đơn qua tin nhắn, ghi tay vào sổ và đối soát chuyển khoản vào cuối ngày: khách không biết món nào
còn, nhân viên phải kiểm kho thủ công, và quán không có nguồn dữ liệu duy nhất cho doanh thu. Bài
toán cốt lõi vì vậy không phải là ``làm một trang bán hàng'', mà là \textbf{giữ cho số lượng bán ra
luôn khớp với nguyên liệu thật trong kho} trong khi khách thanh toán trực tuyến.

\subsection{Mục tiêu}
\begin{itemize}
  \item Số hoá trọn luồng \emph{chọn món $\rightarrow$ tuỳ chọn size/topping $\rightarrow$ đặt trước
        $\rightarrow$ thanh toán $\rightarrow$ nhận tại quầy}.
  \item Quản lý tồn kho theo \textbf{nguyên liệu và công thức} thay vì theo từng sản phẩm, để hai món
        dùng chung một nguyên liệu không thể bán vượt kho.
  \item Thanh toán không dùng tiền mặt qua ZaloPay Sandbox, có chống tạo giao dịch trùng, có đường lui
        khi cổng thanh toán lỗi và giữ phản hồi API dưới 500 ms (NFR-01).
\end{itemize}

\subsection{Phạm vi}
\textbf{Trong phạm vi:} thực đơn và tuỳ chọn, giỏ hàng, đăng ký/đăng nhập, đặt đơn, thanh toán
ZaloPay Sandbox, theo dõi và huỷ đơn, quản trị thực đơn -- kho -- công thức -- đơn hàng -- thanh toán.
\textbf{Ngoài phạm vi:} giao hàng tận nơi, POS tại quầy, đa chi nhánh, hoá đơn điện tử, ứng dụng di
động native và tích hợp máy in bếp.

\section{Tác nhân và mô hình use case}

Hệ thống có năm tác nhân. \textbf{Khách vãng lai} xem thực đơn, tìm kiếm và xem chi tiết món.
\textbf{Khách hàng} đăng nhập, quản lý giỏ, đặt và thanh toán đơn, theo dõi hoặc huỷ đơn.
\textbf{Quản trị viên} quản lý thực đơn, kho, công thức, đơn hàng và thanh toán. \textbf{ZaloPay
Sandbox} là hệ thống ngoài tạo giao dịch và gửi callback xác nhận. \textbf{Cron scheduler} là tiến
trình nền huỷ đơn quá hạn và hoàn lại tồn kho đã giữ.

\begin{figure}[t]
\centering
\includegraphics[width=0.76\textwidth]{../figures/diagrams/BrewLite-UseCase-Overall.png}
\caption{Sơ đồ use case tổng quan, ánh xạ 12 use case tới các yêu cầu chức năng FR-01--FR-11}
\label{fig:usecase}
\end{figure}

\figref{fig:usecase} nhóm 12 use case thành bốn cụm: \emph{khám phá thực đơn} (UC-01--UC-03),
\emph{mua hàng} (UC-04--UC-08), \emph{vận hành đơn} (UC-09) và \emph{quản trị dữ liệu}
(UC-10--UC-12). Ba quan hệ \texttt{include} thể hiện ràng buộc bắt buộc: đặt đơn phải đăng nhập và đã
chọn size, luôn kéo theo giữ tồn kho, và thanh toán là bước nối tiếp. Ba use case chịu rủi ro cao nhất
là UC-06, UC-07 và UC-08, nên được đặc tả chi tiết dưới đây.

\needspace{13\baselineskip}
\section{Đặc tả use case: Đặt đơn hàng và thanh toán (UC-06)}

\begin{table}[htbp]
\centering\small
\caption{Đặc tả use case UC-06}
\label{tab:uc06}
\begin{tabularx}{\textwidth}{@{}>{\raggedright\arraybackslash}p{3.5cm} X@{}}
\toprule
\textbf{Thuộc tính} & \textbf{Nội dung} \\
\midrule
Mã -- Tên & UC-06 -- Đặt đơn hàng và thanh toán trực tuyến \\
Tác nhân chính & Khách hàng đã đăng nhập \\
Tác nhân phụ & ZaloPay Sandbox (cổng thanh toán) \\
Tiên quyết & JWT hợp lệ; giỏ hàng có ít nhất một món; các món còn đủ nguyên liệu. \\
Kích hoạt & Khách bấm \emph{Đặt hàng và thanh toán} tại trang \code{/checkout}. \\
Hậu quyết & Đơn ở trạng thái \code{PENDING} hoặc \code{PAID}; tồn kho đã bị giữ; có bản ghi lịch sử
             trạng thái và bản ghi thanh toán. \\
\bottomrule
\end{tabularx}
\end{table}

\textbf{Luồng chính.}
\begin{enumerate}[label=\arabic*.,leftmargin=1.6em]
  \item Hệ thống xác thực JWT và lấy \code{userId}; trả 401 nếu token hết hạn.
  \item Client gửi danh sách món gồm mã sản phẩm, size, topping và số lượng.
  \item Server đọc bảng giá size và giá topping từ cơ sở dữ liệu, \textbf{tính lại toàn bộ đơn giá và
        tổng tiền}; giá do client gửi lên không được sử dụng.
  \item Hệ thống áp mã khuyến mãi và điểm tích luỹ (nếu có), kiểm tra hạn mức sử dụng điểm.
  \item Với từng nguyên liệu trong công thức của món theo size, hệ thống trừ kho bằng câu lệnh cập
        nhật có điều kiện trên \code{version} (xem Mục~\ref{sec:model}).
  \item Trong một transaction, hệ thống tạo \code{Order}, \code{OrderItem}, các dòng
        \code{OrderIngredient} (ảnh chụp định lượng đã trừ) và \code{OrderStatusHistory}.
  \item Hệ thống tạo bản ghi thanh toán \code{PENDING}, gọi API tạo giao dịch của ZaloPay và trả về
        \code{paymentUrl}; client chuyển sang trang thanh toán.
  \item Khi ZaloPay gọi callback, hệ thống kiểm tra chữ ký HMAC-SHA256 bằng \code{key2} theo cách so
        sánh an toàn thời gian, cập nhật đơn sang \code{PAID} và ghi lịch sử.
\end{enumerate}

\textbf{Luồng phụ và ngoại lệ.}
\begin{itemize}
  \item \textbf{A1 -- Thiếu nguyên liệu:} bước 5 thất bại, toàn bộ transaction được rollback và API
        trả 409 kèm tên nguyên liệu không đủ; không có đơn nào được tạo.
  \item \textbf{A2 -- Cổng thanh toán lỗi hoặc timeout:} đơn giữ nguyên \code{PENDING}; khách được
        phép \emph{đổi phương thức thanh toán} hoặc thử lại trên \textbf{cùng} đơn, không sinh đơn mới.
  \item \textbf{A3 -- Callback đến trễ hoặc không đến:} trang thanh toán tự cập nhật mỗi 15 giây;
        server chủ động truy vấn trạng thái từ ZaloPay; cron huỷ đơn quá hạn và hoàn kho.
  \item \textbf{A4 -- Huỷ đơn:} chỉ cho phép ở \code{PENDING}, \code{PAYMENT\_FAILED} và \code{PAID};
        kho được hoàn theo đúng các dòng \code{OrderIngredient} đã lưu, điểm được trả lại.
  \item \textbf{E1 -- Trùng yêu cầu:} bấm hai lần hoặc callback lặp đều không tạo giao dịch thứ hai,
        nhờ khoá idempotency và ràng buộc duy nhất trên bản ghi thanh toán.
  \item \textbf{E2 -- Hai khách tranh nhau phần nguyên liệu cuối:} một yêu cầu thành công, yêu cầu
        còn lại nhận 409; cơ sở dữ liệu không bao giờ ghi nhận tồn kho âm.
\end{itemize}

\section{User story và tiêu chí chấp nhận}

Mười user story dưới đây bao phủ FR-01--FR-11 và được viết theo góc nhìn người dùng cuối. Cột
\emph{Tiêu chí chấp nhận} chỉ giữ những điều kiện kiểm chứng được, tương ứng với các kiểm thử tự
động trình bày trong báo cáo Assignment~2.

\begingroup\footnotesize
\begin{longtable}{@{}p{0.95cm} p{1.5cm} p{5.0cm} p{6.4cm} p{1.2cm}@{}}
\caption{Danh sách user story và tiêu chí chấp nhận}\label{tab:us}\\
\toprule
\textbf{ID} & \textbf{Vai trò} & \textbf{User story} & \textbf{Tiêu chí chấp nhận} & \textbf{Ưu tiên} \\
\midrule
\endfirsthead
\toprule
\textbf{ID} & \textbf{Vai trò} & \textbf{User story} & \textbf{Tiêu chí chấp nhận} & \textbf{Ưu tiên} \\
\midrule
\endhead
\bottomrule
\endfoot
US-01 & Khách & Tôi muốn xem thực đơn kèm ảnh và giá theo size để chọn nhanh món mình thích. &
Tám món hiển thị đầy đủ; lọc theo danh mục; có trạng thái loading và empty; ảnh có thuộc tính thay thế. & Must \\
US-02 & Khách & Tôi muốn tìm món bằng từ khoá không dấu để không phải gõ đúng chính tả. &
Gõ ``ca phe'' khớp ``Cà phê sữa đá''; không phân biệt hoa thường; lọc ngay khi gõ. & Must \\
US-03 & Khách & Tôi muốn chọn size và topping và thấy giá đổi ngay để biết mình sẽ trả bao nhiêu. &
Size S/M/L đổi đơn giá; topping cộng dồn đúng; tổng tiền khớp với số tiền server tính. & Must \\
US-04 & Khách & Tôi muốn giỏ hàng còn nguyên khi tải lại trang để không phải chọn lại. &
Giỏ lưu cục bộ; sửa và xoá được từng dòng; tổng số ly và tổng tiền cập nhật tức thì. & Must \\
US-05 & Khách & Tôi muốn đăng nhập một lần để đặt đơn mà không phải nhập lại thông tin. &
Đăng nhập trả JWT; sai mật khẩu báo lỗi rõ; vào \code{/checkout} khi chưa đăng nhập thì bị chuyển
tới trang đăng nhập kèm tham số quay lại. & Must \\
US-06 & Khách & Tôi muốn trả tiền online để không cần chuẩn bị tiền mặt. &
Đơn được tạo ở trạng thái chờ thanh toán; chuyển được sang ZaloPay Sandbox; quay lại thì trạng thái
đơn phản ánh đúng kết quả. & Must \\
US-07 & Khách & Tôi muốn theo dõi đơn đã đặt để biết khi nào có thể nhận. &
Bảy trạng thái hiển thị thành dòng thời gian; tự cập nhật mỗi 15 giây; có liên kết tiếp tục thanh
toán khi đơn còn chờ. & Should \\
US-08 & Khách & Tôi muốn huỷ đơn trước khi trả tiền và được hoàn điểm, hoàn kho. &
Huỷ được ở trạng thái cho phép; điểm tích luỹ trở về; kho tăng lại đúng định lượng đã trừ. & Should \\
US-09 & Admin & Tôi muốn đổi trạng thái đơn theo đúng quy trình để tránh nhảy cóc. &
Chỉ chấp nhận chuyển trạng thái hợp lệ; chuyển sai trả 409; mọi thay đổi được ghi vào lịch sử. & Must \\
US-10 & Admin & Tôi muốn quản lý kho, công thức và thực đơn để vận hành hằng ngày. &
Thêm/sửa/xoá nguyên liệu và công thức theo size; cảnh báo khi tồn dưới định mức; nhập kho ghi phiếu. & Must \\
\bottomrule
\end{longtable}
\endgroup

\section{Kiến trúc hệ thống}

\subsection{Kiến trúc tổng quan}
BrewLite là monorepo gồm hai ứng dụng triển khai độc lập, giao tiếp qua HTTPS/JSON: \textbf{ứng dụng
web} (Next.js~15, App Router, React~19, Tailwind, mobile-first) và \textbf{API nghiệp vụ} (NestJS~11,
TypeScript strict). API chia theo miền nghiệp vụ thành chín module: \code{auth}, \code{products},
\code{cart}, \code{orders}, \code{payments}, \code{ingredients}, \code{recipes},
\code{stock-imports}, \code{users}. Tầng dữ liệu là PostgreSQL truy cập qua Prisma~6 với 18 model và
migration chỉ tiến. Mọi endpoint nằm dưới tiền tố \code{/api/v1}; xác thực bằng JWT, phân quyền theo
vai trò ở guard, kiểm tra dữ liệu vào bằng DTO. Phía hạ tầng: web trên Vercel, API trong Docker trên
Render, dữ liệu trên Neon, và GitHub Actions chạy kiểm thử cùng build trên mỗi thay đổi của nhánh
\code{main}.

\subsection{Mô hình dữ liệu kho nguyên liệu và công thức}
\label{sec:model}

Đây là quyết định thiết kế quan trọng nhất của hệ thống. BrewLite \textbf{không lưu tồn kho trên sản
phẩm}. Tồn kho thuộc về \code{Ingredient} (kèm \code{stock}, \code{minStock} và \code{version}),
còn \code{RecipeIngredient} mô tả một sản phẩm ở một size cần bao nhiêu đơn vị của mỗi nguyên liệu.
Khi đặt hàng, hệ thống duyệt công thức để trừ kho; khi huỷ, hệ thống hoàn kho theo ảnh chụp
\code{OrderIngredient} đã lưu.

Lý do: một ly dùng nhiều nguyên liệu và một nguyên liệu phục vụ nhiều món. Cappuccino size S cần
18\,g cà phê hạt Arabica, 120\,ml sữa tươi và 30\,ml kem sữa béo; Americano size S cũng cần đúng
18\,g cà phê Arabica đó. Nếu mỗi sản phẩm giữ một cột tồn kho riêng, hai món sẽ trừ hai cột khác
nhau trong khi kho thật chỉ có một bao cà phê --- và không cột nào phản ánh được việc bán ly này làm
giảm khả năng bán ly kia. Với mô hình công thức, câu hỏi ``còn bán được bao nhiêu ly?'' trả lời được
bằng $\min_i (\text{stock}_i / \text{quantity}_i)$ trên toàn bộ nguyên liệu của món.

\begin{figure}[t]
\centering
\begin{tikzpicture}
  \node[ent] (product) at (-5.45,0)
    {\textbf{Product}\\[1pt]\rule{2.85cm}{0.4pt}\\[1pt]
     id, name\\ categoryId\\ isAvailable};
  \node[ent] (recipe) at (0,0)
    {\textbf{RecipeIngredient}\\[1pt]\rule{2.85cm}{0.4pt}\\[1pt]
     productId, size\\ ingredientId, quantity\\ \emph{(khoá ghép)}};
  \node[ent] (ing) at (5.45,0)
    {\textbf{Ingredient}\\[1pt]\rule{2.85cm}{0.4pt}\\[1pt]
     id, name, unit\\ stock, minStock\\ version};

  \draw[lnk] (product.east) -- node[above,card]{1} node[below,card]{N} (recipe.west);
  \draw[lnk] (recipe.east) -- node[above,card]{N} node[below,card]{1} (ing.west);

  \node[note,text width=13.4cm,anchor=north] at (0,-1.35)
    {Giá theo size nằm ở \texttt{ProductSizePrice}; khi đặt hàng, mỗi dòng công thức sinh một dòng
     \texttt{OrderIngredient} lưu ảnh chụp định lượng đã trừ.};
\end{tikzpicture}
\caption{Lược đồ rút gọn: sản phẩm -- công thức -- nguyên liệu}
\label{fig:erd}
\end{figure}

Vì tồn kho là tài nguyên dùng chung, thao tác trừ kho phải an toàn khi nhiều đơn đến cùng lúc.
BrewLite dùng \emph{optimistic locking}: mỗi nguyên liệu có cột \code{version}, và lệnh cập nhật chỉ
khớp khi \code{version} chưa đổi \textbf{và} \code{stock} vẫn đủ cho lượng cần trừ; khi đó hệ thống
giảm \code{stock}, tăng \code{version} và kiểm tra số dòng bị ảnh hưởng. Nếu hai transaction cùng đọc
một \code{version}, chỉ một lệnh khớp điều kiện; lệnh còn lại nhận \code{count = 0} và bị biến thành
lỗi 409 thay vì ghi đè. Toàn bộ việc tạo đơn chạy trong transaction mức \code{Serializable}, nên đơn,
chi tiết đơn, ảnh chụp nguyên liệu và lịch sử trạng thái hoặc cùng tồn tại, hoặc cùng biến mất.

\subsection{Các quyết định thiết kế then chốt}
\begin{itemize}
  \item \textbf{Giá tính ở server, định lượng lưu bằng ảnh chụp.} Client chỉ gửi mã món, size, topping
        và số lượng; đơn giá lấy từ cơ sở dữ liệu nên sửa giá trên trình duyệt không làm đổi số tiền
        phải trả. \code{OrderIngredient} lưu lại đúng lượng đã trừ, nhờ vậy nếu công thức thay đổi sau
        đó thì việc huỷ đơn vẫn hoàn kho đúng như lúc đặt.
  \item \textbf{Máy trạng thái đơn hữu hạn.} Bảy trạng thái (\code{PENDING}, \code{PAYMENT\_FAILED},
        \code{PAID}, \code{PREPARING}, \code{READY}, \code{COMPLETED}, \code{CANCELLED}) với bảng
        chuyển đổi tường minh; mọi bước chuyển đều ghi lịch sử.
  \item \textbf{Thanh toán idempotent và biên giới tin cậy.} Mỗi giao dịch gắn với một mã duy nhất
        và một ràng buộc duy nhất trong cơ sở dữ liệu, nên callback lặp lại hay bấm hai lần đều không
        tạo giao dịch thứ hai. Dữ liệu từ ZaloPay chỉ được chấp nhận sau khi xác thực chữ ký; khoá bí
        mật chỉ nằm trong biến môi trường, không có trong mã nguồn hay log.
\end{itemize}

\section{Kết luận}
Báo cáo đã xác lập 12 use case, 10 user story kèm tiêu chí chấp nhận và kiến trúc ba lớp. Điểm khác
biệt nằm ở mô hình dữ liệu: tồn kho gắn với \textbf{nguyên liệu và công thức} và được bảo vệ bằng
optimistic locking cùng transaction khả tuần tự, nên hai món dùng chung nguyên liệu không thể cùng
bán vượt kho. Đây là nền cho các bước tiếp theo --- đo hiệu năng, kiểm thử tự động toàn bộ luồng đặt,
thanh toán, huỷ đơn và hoàn thiện triển khai --- trình bày trong báo cáo Assignment~2.

\end{document}


